\subsection{Overview}

The literature on agricultural subsurface (tile) drainage repeatedly converges on the same meta-conclusion: many of the most policy-relevant questions about nutrient loss, water balance, and soil evolution in drained landscapes---and especially in drained organic and problem soils---cannot be resolved with the short, single-site, single-nutrient datasets that dominate the field. Across a wide geographic span, authors call for (i) \emph{longer observation periods} that span climatic variability and slow soil evolution, (ii) \emph{better metrics}, particularly phosphorus (P) saturation and sorption indices that outperform standard agronomic soil tests, (iii) \emph{isotope- and tracer-based approaches} to resolve loss pathways and source attribution, (iv) \emph{multi-scale and multi-species designs} that partition water and nutrient loads among competing pathways, and (v) \emph{mechanistic and modelling advances} that couple hydrology, biogeochemistry, and human-induced soil change. \textcite{Gramlich2018} frame the headline gap directly: longer-term studies---especially those exceeding three to four years---on drained organic soils are needed with respect to water balance and nitrogen (N) losses, along with case studies that compare drained and undrained conditions directly for N, P, and plant-protection products. This section organises the field's open questions and stated priorities around these lines of evidence, emphasising what remains uncertain, which hypotheses need testing, and which methods and data are still missing.

\subsection{The recurring call for long-term, multi-climate, multi-site observation}

The single most common recommendation is simply to keep observing for longer and across more conditions. \textcite{Kokulan2019} argue that additional longer-term field studies are needed to capture climatic variability and that continued observation is required to evaluate whether and how field hydrology and soil structure evolve into the future. Similar multi-year, multi-climate appeals recur for N-fertiliser effects on yield \parencite{Kaur2024}, for high-carbon soil mineralization and management \parencite{Waring2022}, and for developing management strategies that balance economic production with environmental protection across soil types and climatic zones \parencite{Kladivko2004}. \textcite{RANDALL1995} stress that long-term studies conducted over a range of climatic conditions are essential to derive appropriate guidelines, and that detailed plant, soil, and water measurements must be gathered together if interactions among crop production, climate, and environmental impact are to be uncovered.

Long-term designs are also invoked to detect slow or subtle change. \textcite{Kumar2014} note that because tillage and drainage differences in their system were small, long-term monitoring of crop yield is important, and that a switch from continuous corn to a corn--soybean rotation may help regain soil organic carbon over long durations. \textcite{Manninen2023} recommend keeping long-term field studies running to obtain more accurate estimates of discharge-induced organic-matter export from various soil types and cultivation practices, and to clarify the processes behind soil organic carbon and N losses. \textcite{Frison2009} extend the horizon further, asking whether observed drainage-induced soil evolution will continue or stabilise over decades and what the consequences will be for agricultural and environmental soil function; they also call for measurements along drains at different slope positions and on drainage networks of different ages. Long-term operation is itself framed as a management imperative: \textcite{Singh2007} argue that fields with subsurface drainage must continue to operate the system and keep reclaimed land under cultivation to maintain soil health on a long-term basis.

\subsubsection{Sampling standardisation and consistency}

Several authors emphasise that longer records are only useful if sampling is standardised. \textcite{Sherbine2023} recommend that soil biological indicators be evaluated over longer periods and that sampling be done at consistent times each year under similar soil moisture and temperature conditions. \textcite{Li2023b} likewise argue that the timing of sampling for soil health index (SHI) determination should be standardised to enable comparable results, practically aligning with post-harvest sampling for standard soil nutrient testing. For meta-analytic comparability, \textcite{Ghirardini2019} call for a complete hydrologic description of each plot, standardised sampling and analytical methods, monitoring periods longer than 40 days (especially for dewatered sludge), and reporting of concentrations with descriptive statistics (minimum, maximum, standard deviation).

\subsection{Phosphorus: better metrics, pathways, and loss prediction}

Phosphorus dynamics attract the densest set of open questions, spanning the adequacy of soil tests, the mechanisms of desorption and transport, and the tools needed for prediction.

\subsubsection{Beyond agronomic soil-test P: saturation and sorption metrics}

A consistent theme is that agronomic soil-test P (STP) is insufficient for environmental risk. \textcite{Sharpley1995} argue that agronomic P tests may not reliably assess eutrophication potential and that environmental tests estimating bioavailable P and P sorption capacity should be used on soils flagged as ``high risk,'' while comprehensive P-management plans address soils already containing P in excess of crop requirements. \textcite{Hooda1999} similarly propose that a measure of the degree of soil P saturation (DPS) may better estimate loss potential than extractable P such as Olsen-P, and that a full accounting of all P output forms is needed for mitigation strategies. \textcite{Nguyen2002} call for establishing critical DPS values---including for New Zealand sediments---calibrated across a range of sediment P tests and equilibrium P concentrations at net-zero sorption--desorption. \textcite{MILLER1979} adds a design-stage caution: P-adsorption capacity should be considered before organic soil areas are developed for crop production. \textcite{Osterholz2020} recommend assessing additional fields in the 100--400~mg~kg\textsuperscript{-1} STP range to clarify STP--P relationships and examining more detailed soil properties, including P sorption saturation, macropore formation and stability, and mineralogy.

\subsubsection{Mechanisms of desorption, colloidal and particulate transport}

Multiple authors identify the sorption/desorption interface and particulate pathways as poorly resolved. \textcite{Mumbi2025} call for quantifying when and how varying water sources and fluxes interact with soil to desorb and adsorb P---a need considered critical to advancing prediction of subsurface dissolved reactive P (DRP) transport---supported by data across soil types and management (e.g., long-term till vs.\ no-till). \textcite{Flores2021} note the need to understand the mechanisms controlling particulate settling and P desorption. \textcite{Nguyen2002} call for investigating P release--retention characteristics of drainage sediments under prolonged flooding and alternating waterlogged--drying conditions, and the interaction of sediment P with overlying drainage water across flooding, drying, wetting, and P-loading regimes. \textcite{Jiang2010} recommend assessing P fractions after leaching to understand transformations under polyacrylamide (PAM), including its effect on different purple-soil layers and interactions among broadcasting rate, mixing procedure, and leaching losses. \textcite{Wilson2019} highlight the need to quantify the contribution of colloidal P to DRP measurements and to determine whether erosion-mitigating BMPs or nutrient management better control it, alongside clarifying P stratification, preferential flow through tile drainage, and the effectiveness and time-to-effect of subsurface P application in stratified fields. \textcite{Norgaard2013} recommend including both dissolved and colloid-associated leaching parameters in field-scale risk assessment, testing a modified Dexter concept using clay and silt as a linking indicator, and considering water-dispersible colloids and tillage events as indicators of particle release.

\subsubsection{Multi-species reporting, temporal effects, and databases}

Fragmented reporting hampers synthesis. \textcite{Christianson2016} call for increasing the number of field-scale studies documenting drainage P losses and the effects of cropping, nutrient application, soil properties, and drainage design; including multiple P species (DRP, particulate P, total P); further developing the Drain Load database to capture peak flowrate and a drainage-concentration table; and using individual-study source data to evaluate temporal effects on drainage P transport. \textcite{Macrae2019} call for examining the hydroclimatic drivers of event responses and their impact on differing P concentrations in tile drainage, determining how tillage or P placement affects such relationships, and improving understanding of small-scale hydrologic and biogeochemical processes under various flow conditions and soil types. \textcite{Tanga2025} note that controlled-drainage effects on soluble P remain variable in the literature and recommend monitoring both particulate and dissolved forms while selecting site-specific stoplog depths based on field slope, drain slope, soil type, and climate. \textcite{Grant1996} call for further study of soil properties and P transport and for determining the applicability of in-situ optical turbidimeters to estimate particulate-matter content in drainage water. \textcite{Daigh2015b} note that subsurface-drainage total reactive P concentrations and yields in bioenergy systems need further evaluation in areas prone to higher P losses.

\subsubsection{Organic (Histosol/peat) soil phosphorus}

Phosphorus in cultivated organic soils is singled out as poorly understood. \textcite{Grenon2022} call for mechanistic studies to elucidate the pathways, hot spots, and hot moments of P transfer from agricultural organic soils to adjacent water bodies, including the dynamics of residual P fractions, soil--water P dynamics through long-term water-quality study, and P-sorption dynamics to compute more accurate P balances. \textcite{Grenon2021} propose a concrete methodological agenda: assessing P geochemistry in agricultural Histosols using isotope tracing (\textsuperscript{32}P/\textsuperscript{33}P and oxygen isotope ratios in phosphates), XANES, and \textsuperscript{31}P~NMR alongside sequential fractionation; assessing microbial turnover times; conducting long-term studies of drainage water management to reduce total P loads; developing a new mathematical model accounting for geo-mechanical instability and variability in organic-soil properties; and creating P-transport indices to predict loss. \textcite{Höyhtyä2025} note that nutrient leaching from abandoned peat fields is poorly studied, calling for research into the effects of ditch blockage, energy-willow production, or forestation, and for exploring the true nature of H\textsubscript{2}SO\textsubscript{4}-P via P K-edge XANES spectroscopy.

\subsubsection{Isotope-exchange and kinetic prediction of P loss}

\textcite{McDowell2001} propose that P isotope-exchange kinetics (IEK) has the potential to explain and eventually predict P loss to overland flow and subsurface drainage across a wide range of STP concentrations, and that with further work the IEK approach may predict P loss more accurately.

\subsection{Nitrogen: mineralization, testing, and management options}

Nitrogen questions cluster around residual and mineralizable N, testing methods, and the search for management levers beyond fertiliser rate. \textcite{Waring2022} call for measuring mineralization rates from high-carbon soils and for N-management options beyond fertilisation practices to reach water-quality reduction goals. \textcite{Bosompemaa2021} recommend further quantifying the mass of N in soil during both growing and non-growing seasons. \textcite{RANDALL1995} suggest a soil N test in the spring following dry years to measure residual nitrate and credit it against fertiliser applications, while \textcite{Conrad2009} recommend measuring mineral N to evaluate leaching potential before a new crop is sown and modifying crop rotation to prevent long periods without soil cover or N uptake. \textcite{Gentry2014} argue that more accurate quantification of fertiliser-N rates and timing would sharpen estimates of fertiliser-management influence on nitrate transport, and that the appropriateness of fall N application on tile-drained land in east-central Illinois may need revisiting in light of warmer winters. \textcite{ElGhannam2023} recommend that future studies measure not only nitrate outflow but also soil nitrate and plant N status. \textcite{Saiki2020} call for assessing the effect of water flow on soil N accumulation and discharge by evaluating N balance in paddy fields. On tooling, \textcite{Masri2024} suggest resin-membrane strip technology as a future tool to characterise soil N cycling and release under variable conditions, and \textcite{Kaur2024b} point to technology that improves the efficiency of enhanced fertiliser management in tile-drained fields. \textcite{Aide2016} frame an integrated agenda: developing effective crop-nutrient-management systems to improve uptake efficiency and reduce nitrate leaching, and refining denitrification-bioreactor design (reactor size, bed packing materials, equilibrium thresholds, elimination of preferential flow paths, and pH maintenance).

\subsubsection{Predicting long-term nitrate concentrations}

\textcite{Skaggs2012} identify a distinct methodological gap in estimating long-term flow-weighted average N concentrations for drainage waters: methods are needed for benchmark soils under typical rotations; the approximate method requires improvement for years when concentrations deviate strongly from the long-term average; and additional testing across soils, crops, and climates is needed, ideally with simple methods relating concentrations to prior yields, fertilisation, and weather patterns.

\subsection{Isotopes, tracers, and pathway partitioning}

A cross-cutting methodological priority is the use of isotopic and tracer techniques to resolve otherwise ambiguous pathways and sources. \textcite{Grenon2022} argue that more frequent post-fertilisation sampling and isotope tracers would help elucidate P dynamics in cultivated organic soils, and that applying isotopes within soil--water dynamic studies could shed new light on loss pathways. \textcite{Islam2022} propose measuring the stable isotope composition of water ($\delta$D and $\delta^{18}$O) to identify plant water sources late in the season, and \textcite{YliHalla2022} recommend verifying preliminary groundwater flow-path estimates with stable water isotopes. \textcite{Komor1997} show a source-attribution application, calling for additional $\delta^{11}$B analyses from closely spaced wells around feedlots and cultivated fields to define spatial variation in manure-leachate proportions, and for further investigation of the isotopic compositions of B sources to strengthen $\delta^{11}$B as a contaminant tracer. At the pathway-partitioning level, \textcite{Flores2021} urge future drainage studies to comment explicitly on surface-inlet presence and to partition both water and nutrient loads by pathway using before-after control-impact designs, hydrograph separation, or tracers, and to report both N and P loads for each site-year (including additional N forms such as NH\textsubscript{4} and organic N). \textcite{Kokulan2019} similarly recommend evaluating the chemistry and relative contributions of overland flow (OF) versus tile flow (TF) to edge-of-field nutrient losses. \textcite{Laubel1999} advocate combining studies at different scales---intact soil columns, experimental plots, and catchments---to examine subsurface loss of particulate matter and P in relation to macroporosity, soil P content, and cultivation practice.

\subsection{Soil biology, health indicators, and human-induced soil evolution}

\subsubsection{Biological and soil-health indicators}

The soil-biology response to changing drainage hydrology is repeatedly flagged as understudied. \textcite{Sherbine2023} call for longer-term evaluation of soil biological indicators to understand how soil biology responds to gradually changing hydrology. \textcite{Yang2025} identify the interaction mechanisms between subsurface drainage (SSD) and soil physicochemical properties, and the response of soil microbial activity, as priorities. \textcite{Kaur2024b} and \textcite{Kumar2014} call for continued research on specific soil microbial properties and microbial community structure and function to explain tillage and drainage impacts. \textcite{Li2023b} identify a scalability gap for nematode measurements, recommending DNA-sequencing methods and statistical approaches that quantitatively integrate nematodes with the SHI, and provide a basis for expanding to multiyear assessments of SHI linkages with nutrient losses and crop productivity. \textcite{Nakajima2015} argue that additional long-term studies are needed on the effects of no-till and subsurface drainage on the soil quality index (SQI), that valid, reliable, sensitive, repeatable, and accessible indicators must be identified, and that a framework for overall soil-quality evaluation and a systems approach are needed. \textcite{Singh2007} suggest that soil-enzyme activities could help identify management practices to maintain soil health. \textcite{Masri2024} argue that further research is needed to prove that no-till improves soil-health functions enough to sustainably reduce nitrate leaching, and that an extended research period is needed to relate nitrate load, water-quality parameters, and soil-health functions. \textcite{Caron2025} recommend adding dynamic to static indicators to diagnose compaction (extending to deeper layers and more corn-producing areas) and investigating whether long-term forage rotation or agroforestry can restore soil physical health, including comparative long-term economic viability.

\subsubsection{Drainage-induced soil evolution over decadal scales}

A distinctive line of inquiry treats subsurface drainage as a driver of long-term, potentially chaotic soil evolution. \textcite{Montagne2009} argue that quantifying human-induced soil evolution as a function of time should be a research priority, that further research is needed to evaluate drainage's impact on clay-mineral dissolution, and that more quantitative data as a function of soil properties are needed to model impacts on soil functions. \textcite{Montagne2010} argue that soil evolution in response to climate change must be explicitly studied and included in global climate-change models. \textcite{Montagne2013} suggest that significant material losses over a few decades make human-induced perturbation a powerful tool for studying the time-scales of such chaotic soil evolution. \textcite{HOOGHOUDT1952} note the continuation of a long experiment to determine whether soil structure will eventually be affected. \textcite{Ouellet2008} recommend using variables expressing tillage practice, surface-residue coverage, years in no-till, and soil particle size for empirical prediction (considering variable interactions), suggest soil organic matter as a crude surrogate for residue cover in regional modelling, and note that predictions of earthworm activity would likely be stronger for no-till than conventional-tillage systems.

\subsection{Drainage water management, BMPs, and policy tools}

Several authors target the effectiveness and siting of best management practices (BMPs) and drainage water management (DWM). \textcite{Wilson2019} present an extensive agenda: determining the impact of DWM on surface and tile flow of water and P; mapping the spatial distribution and range of STP values (especially elevated fields); improving understanding of P cycling at in-field, edge-of-field, and in-stream scales; developing innovative policy solutions (manure cooperatives, transport incentives, market-based strategies) and assessing farmer willingness and site-specific cost-benefit relationships; creating standardised methods to track adoption; and developing easy-to-use decision-support tools and field-scale P-loss projections to guide BMP placement. \textcite{Singh2020c} likewise call for extensive long-term, field- and watershed-scale work: measuring BMP effectiveness in site-specific areas; long-term field-scale trials of BMPs serving as P sinks; studies of adsorption/desorption with P-binding amendments; investigation of the time and pathways over which flushes of available P are lost or fixed; simulation of measured P loss from bioreactors and saturated buffers supplemented with P-adsorption materials; field evaluation of the P-index and critical-source-area concepts; long-term studies on P stratification linked to water-quality impairment; field-scale evaluation of lime, gypsum, and slow-release P fertilisers; further work on P source, rate, placement, and timing; year-round natural-rainfall runoff studies; multi-year cover-crop and tillage studies; long-term monitoring of new controlled-drainage, bioreactor, and saturated-buffer designs; better-funded watershed-scale P-flux monitoring; and criteria to maintain and remove P from BMPs before retention capacity is exceeded. \textcite{Zhang2017} add that impacts of modified no-till on DRP losses need investigation and that new soil-P testing and optimised P application timing, rates, and methods are needed to reduce DRP losses under no-till. \textcite{Nash2015c} call for determining whether an observed P-loss response between managed and free drainage was due to soil type, pH, anaerobic conditions, shallow tile placement, plant uptake, and/or rainfall. \textcite{Gottschall2012} argue that field-scale studies of environmental concentrations and loading across different soil-management systems are needed to validate predicted concentrations, elucidate dominant exposure pathways, and characterise dissipation kinetics.

\subsection{Modelling, reactive transport, and specialised soils}

\subsubsection{Model development and reactive transport}

Modelling gaps span parameterisation, coupling, and validation. \textcite{Conrad2009} recommend a range of parameter values rather than single estimates, improved understanding of organic--inorganic transformations in soil and plant (including plant growth), and statistically based automatic calibration. \textcite{Kokulan2019} note that runoff models for tiled landscapes with cracking clay soils can be improved by including temporal variability in the activation of hydrological pathways and hydroclimatic drivers of activation patterns. \textcite{Dadfar2010} plan to incorporate clay-mineralogy data into a CF algorithm as it becomes available, to better define the clay-content threshold, and to complete field validation against water-quality data from several Canadian watersheds. \textcite{KARAJEH1994c} propose extending a reactive solute-transport formulation to account for lime (CaCO\textsubscript{3}) precipitation/dissolution as a salt sink/source and for the reactivity and mobility of toxic trace elements (selenium, boron, molybdenum, arsenic) to serve irrigation, drainage, and shallow-groundwater management where salinity and trace elements matter. \textcite{Singh2020c} also call for simulating and validating measured P loss from amendment-supplemented bioreactors and saturated buffers.

\subsubsection{Acid sulfate and boreal soils}

Drained acid sulfate (AS) soils raise specialised questions. \textcite{Salo2023} recommend that future model applications consider short-term hydro-meteorological effects on sulfide oxidation under different drainage setups, and that soil macroporosity, the effect of ripening on hydraulic properties, formation of secondary minerals and weathering, and nitrate and gaseous emissions be incorporated, with the effects of moisture, temperature, and soil structure on sulfide oxidation examined. \textcite{Virtanen2017} note that the role of the N cycle in sulfide oxidation is unknown in boreal cultivated AS soils and that sulfide oxidation coupled with nitrate reduction, together with alternative neutralising agents, warrants study. \textcite{YliHalla2022} call for detailed study of the peat-cover thickness required to hinder sulfidic-material oxidation. \textcite{Åström2005} call for more detailed maps of actual and potential AS soils, more knowledge of the relative contribution of anthropogenic versus organically derived acidity in forested catchments, new environmentally friendly ways of draining actual AS soils, and cessation of drainage of potential AS soils.

\subsubsection{Saline, waterlogged, and reclaimed soils}

For reclaimed saline and waterlogged Vertisols, \textcite{Mukhopadhyay2023} and \textcite{Mukhopadhyay2023b} call for studies of how SSD integrated with crop management (inorganic fertilisers, crop residues, compost, green manure) affects soil-quality improvement and carbon-sequestration potential, of N/P/K losses and field-scale budgeting across more sites and larger Vertisol areas, of efficient recycling of drained water, and of greenhouse-gas fluxes before and after SSD installation to assess climate-mitigation potential. \textcite{ElGhannam2023} call for long-term study of controlled drainage's effect on soil salinity and for site-specific salinity-management practices to avoid salinization under long-term shallow groundwater. \textcite{Colombani2016} call for investigation of salt-crust formation and the subsequent transport of dissolved species, detailed soil-physical analysis of long-term tillage, and more sustainable practices (accurate fertiliser scheduling, frequent small irrigations after fertilisation, and buffer zones to abate excess N). \textcite{SoltaniGerdefaramarzi2020} suggest investigating leaching effects on other elements such as plant-available potassium, analysing at least three soil samples at longitudinal distances from the drain and at different depths by combined sampling, and testing organic modifiers such as animal manures to improve leaching efficiency. \textcite{Qian2021} recommend monitoring soil properties such as elemental toxicities in each year of multi-year waterlogging experiments, extending research to different crops and regions, and scheduling field drainage to account for different waterlogging forms, especially during budding, flowering, and boll-forming stages. \textcite{Bahçeci2018} recommend further field research to verify long-term resistance to root growth (especially for perennial crops) and performance under other soil, hydrological, and agricultural conditions, while \textcite{Bahceci2008} note that drainable-porosity ($\mu$) values can be determined in the laboratory, in situ, or by correlation, with more reliable values obtained when soil texture is considered. \textcite{Sadhukhan2019} call for further tests at additional locations under different fertilisation (including manure), soil, climate, and crop conditions over longer periods with more observational data. \textcite{Isa2023} call for studies of additive effectiveness (soil strength, rate, cost, impact on rice, physicochemical properties, and soil interaction), land-replacement characterisation and economic evaluation, and integration of rehabilitation methods to maximise soil strength and paddy yield.

\subsection{Other emerging topics and cross-cutting data needs}

Several findings point to under-studied constituents and processes. \textcite{Sethna2022} call for analysis of particulate silicon (amorphous and biogenic) to quantify total-Si transport and cycling, study of how warming and changing precipitation affect phytolith mineralization and dissolved-Si inputs to freshwaters, and explicit studies of nutrient limitation and uptake to quantify the role of diatoms and ambient nutrient stoichiometry. \textcite{Manninen2023} recommend that future discharge monitoring focus on variations in nutrient stoichiometry to assess the fate of organic matter in aquatic ecosystems. \textcite{Ghirardini2019} call for addressing the release of accumulated compounds during intense rain and drought/wet cycles, microbial indicators in surface runoff and tile drainage across rainfall events and soil types, and the occurrence of antibiotic-resistant genes and bacteria in sludge-amended soils, alongside comparative studies of sludge types, application rates and methods, and irrigation frequencies on different soils. \textcite{Gramlich2018} additionally call for clarifying the relevance of preferential flow in loamy and sandy soils and for defining categories of relevant compounds and soil--site characteristic combinations that create high cumulative environmental risk under intensive agriculture. \textcite{DambenieceMigliniece2018} recommend detailed fertiliser-use records to confirm crop-area/nutrient-load correlations and analysis of seasonal changes in nutrient concentrations and loads. \textcite{Norberg2025} note that long-term maintenance of soil fertility will probably require additional P and K fertilisation to preserve availability. \textcite{Kladivko2004} argue that detailed analysis of individual storm hydrographs is needed to test event-scale hypotheses and that comparing low-organic-matter soils with the high-organic-matter soils of much of the U.S.\ Midwest is essential for understanding the system and designing region-appropriate management. Collectively, these priorities define a research frontier in which longer records, richer metrics, isotopic tracing, and multi-scale, multi-species monitoring are prerequisites for the mechanistic understanding and predictive models the field still lacks.